problem:  
Let \(W_{p,q} = \mathrm{SU}(3)/\mathrm{U}(1)_{p,q}\) denote an Aloff–Wallach 7‑manifold.  
It is known that for special integer pairs \((p,q)\), the manifold \(W_{p,q}\) admits **homogeneous nearly parallel \(G_2\)** structures \(\varphi_{p,q}\) satisfying  
\[
d\varphi_{p,q} = \lambda *\varphi_{p,q},
\]
for some constant \(\lambda \neq 0\); these are among the finitely many homogeneous nearly parallel \(G_2\) manifolds classified by Friedrich–Kath–Moroianu–Semmelmann.  

Open question:  
For a fixed pair \((p,q)\) that yields a homogeneous nearly parallel \(G_2\) structure \(\varphi_{p,q}\), do there exist **non‑homogeneous deformations** \(\varphi'\) of \(\varphi_{p,q}\) satisfying  
\[
d\varphi' = \lambda' *\varphi', \quad \lambda' \in \mathbb{R}\setminus\{0\},
\]
with \(\varphi'\) not equivalent to \(\varphi_{p,q}\) by diffeomorphism or overall scaling?  
Equivalently, is the infinitesimal deformation space of the nearly parallel \(G_2\) structure on \(W_{p,q}\) nontrivial beyond the \(\mathrm{SU}(3)\)-invariant directions, or are all nearly parallel \(G_2\) structures on each \(W_{p,q}\) rigid and homogeneous?  

Why is it a "good" problem:  
This formulation is both mathematically correct and deeply open—it isolates a precise geometric and analytic rigidity question about the Aloff–Wallach spaces, sitting at the confluence of special holonomy theory, Lie group symmetry, and geometric analysis.

1. **Logical and conceptual soundness:**  
   The problem is correctly phrased in the intrinsic terms of nearly parallel \(G_2\) geometry. It asks about the moduli of solutions to \(d\varphi = \lambda *\varphi\) on a fixed compact 7‑manifold, a well-defined and coherent analytic framework directly tied to special holonomy via cone constructions.

2. **True research-level openness:**  
   While homogeneous nearly parallel \(G_2\) manifolds are fully classified, it is unknown whether non‑homogeneous nearly parallel solutions exist on the same topological spaces. No general theorem rules out such deformations, and the deformation complex for nearly parallel \(G_2\) structures is subtle, governed by the spectrum of the Laplacian acting on coclosed 3‑forms. Determining rigidity or flexibility in this setting is an open and active question in differential geometry.

3. **Analytic–algebraic bridge:**  
   The Aloff–Wallach family brings together **Lie-theoretic symmetry** (via the homogeneous presentation) and **elliptic PDE analysis** (through the nearly parallel condition). Studying whether the homogeneous nearly parallel forms exhaust all solutions requires tools from both representation theory (to decompose invariant tensors) and geometric analysis (to study eigenforms beyond the invariant subspace).

4. **Topological and geometric significance:**  
   The rigidity or non-rigidity of nearly parallel \(G_2\) structures directly influences the landscape of asymptotically conical \(\mathrm{Spin}(7)\) geometries, as every such structure defines the link of a potential Spin(7) cone. A discovery of non‑homogeneous nearly parallel forms would expand known families of exceptional holonomy metrics and shed new light on moduli space phenomena.

5. **Simplicity and depth:**  
   The central question—“Are all nearly parallel \(G_2\) structures on Aloff–Wallach spaces homogeneous?”—is concise and accessible, yet solving it requires mastery of representation theory, curvature identities, and sophisticated elliptic theory. It thus embodies the ideal of a **simple but profoundly deep** research problem.

In summary, this problem cleanly formulates a high‑level open question at the interface of homogeneous geometry and special holonomy analysis, making it a genuinely **good** research‑level problem.